\section{The term}
\label{sec:6.1}
Every other term in the floor governs how an act is done. The first governs whether to take part at all: don't take part in the use of armed force. It comes from Article 2(4) of the Charter, which prohibits the threat or use of force against the territorial integrity or political independence of any state. It has the Charter's two exceptions, self-defense against an armed attack (Article 51) and force the Security Council authorizes under Chapter VII within the scope it authorizes, and two of my own, for stopping a genocide under way and for ending an occupation a disinterested body has found unlawful. In short: \emph{corroborated defense}, \emph{authorized force}, \emph{the genocide exception} and \emph{the occupation exception}.

A Council authorization counts only for the scope it authorizes. A wider reading is the claim of the state acting under the authorization, and is discounted as one.

This term rests on the same foundation as the others, applied between states. A use of force subjects a population to a decision it had no standing to contest, and the claim that licenses it, that the force is defensive, is the belligerent's own. The requirement of corroboration is the own-case rule applied between states.

The term comes in two versions, and every case in this chapter says which it runs on. The first, \emph{the Charter as courts read it}, takes the clause as courts and tribunals read it, with no departure in force: Article 2(4), Article 51 on the defender's own claim until the Council acts, and Chapter VII, together with the Covenant's Article 6, under which a state must protect the lives of people within its jurisdiction from attack. Under Article 51 or Article 6 alike, the licence to repel an attack covers only the least force that repels it, and ends when the attack ends. The second, \emph{the extended term}, adds the departures this chapter argues for: corroborated defense, and standing to defend and the genocide and occupation exceptions. The first rests on law as it stands; the second adds my argument to it. Ukraine, Iran and the domestic clause run on the Charter as courts read it. Gaza runs on both, and the Gaza section labels each step. Taiwan needs standing to defend and corroborated defense, and Cyprus and Karabakh the occupation exception; a later section states those departures, and Appendix~\ref{sec:F} works the cases out.
